\documentclass[12pt,oneside]{book}

% ==========================================================
% METADATOS
% ==========================================================
\newcommand{\universidad}{Universidad de Buenos Aires (UBA)}
\newcommand{\facultad}{Facultad de Derecho}
\newcommand{\carrera}{Carrera de Abogacía}
\newcommand{\materia}{Derecho Procesal Civil}
\newcommand{\titulo}{Proceso de desalojo: la audiencia preliminar (art. 360 CPCCN) y taller de role playing}
\newcommand{\autor}{Isaac Edgar Camacho Ocampo}
\newcommand{\anio}{2026}

% ==========================================================
% PAQUETES
% ==========================================================
\usepackage[utf8]{inputenc}
\usepackage[T1]{fontenc}
\usepackage[spanish,es-nodecimaldot]{babel}

\usepackage[
    top=2.5cm,
    bottom=2.5cm,
    left=3cm,
    right=2.5cm,
    headheight=15pt
]{geometry}

\usepackage{amsmath}
\usepackage{amssymb}
\usepackage{mathtools}

\usepackage{graphicx}
\usepackage{float}
\usepackage{caption}
\usepackage{subcaption}

\usepackage{booktabs}
\usepackage{array}
\usepackage{longtable}
\usepackage{tabularx}

\usepackage{enumitem}

\usepackage{xcolor}
\usepackage[most]{tcolorbox}

\usepackage{fancyhdr}
\usepackage{ragged2e}

\usepackage[
    colorlinks=true,
    linkcolor=blue!50!black,
    citecolor=green!40!black,
    urlcolor=blue!60!black,
    pdftitle={\titulo},
    pdfauthor={\autor},
    pdfsubject={Apuntes universitarios},
    bookmarks=true
]{hyperref}

% ==========================================================
% CONFIGURACIÓN
% ==========================================================
\graphicspath{
    {./img/}
    {./graficos/}
    {./figuras/}
}

\setcounter{tocdepth}{2}

\setlength{\parindent}{0pt}
\setlength{\parskip}{0.5em}

\setlist[itemize]{
    topsep=0.3em,
    itemsep=0.2em
}

\setlist[enumerate]{
    topsep=0.3em,
    itemsep=0.2em
}

\clubpenalty=10000
\widowpenalty=10000

\pagestyle{fancy}
\fancyhf{}
\fancyhead[L]{\nouppercase{\leftmark}}
\fancyhead[R]{\materia}
\fancyfoot[C]{\thepage}
\renewcommand{\headrulewidth}{0.4pt}
\renewcommand{\footrulewidth}{0pt}

\fancypagestyle{plain}{
    \fancyhf{}
    \fancyfoot[C]{\thepage}
    \renewcommand{\headrulewidth}{0pt}
}

% ==========================================================
% COMANDOS
% ==========================================================
\newcommand{\abs}[1]{\left|#1\right|}
\newcommand{\norm}[1]{\left\lVert#1\right\rVert}
\newcommand{\vect}[1]{\mathbf{#1}}
\newcommand{\deriv}[2]{\frac{d#1}{d#2}}
\newcommand{\pderiv}[2]{\frac{\partial#1}{\partial#2}}

% ==========================================================
% ENTORNOS / CAJAS
% ==========================================================
\newtcolorbox{definition}[1]{
    enhanced,
    breakable,
    title={Definición: #1},
    fonttitle=\bfseries,
    colback=blue!3,
    colframe=blue!50!black,
    boxrule=0.8pt,
    arc=2mm
}

\newtcolorbox{important}{
    enhanced,
    breakable,
    title={Importante},
    fonttitle=\bfseries,
    colback=yellow!8,
    colframe=orange!70!black,
    boxrule=0.8pt,
    arc=2mm
}

\newtcolorbox{example}[1]{
    enhanced,
    breakable,
    title={Ejemplo: #1},
    fonttitle=\bfseries,
    colback=green!4,
    colframe=green!40!black,
    boxrule=0.8pt,
    arc=2mm
}

\newtcolorbox{formula}[1]{
    enhanced,
    breakable,
    title={Fórmula / Regla: #1},
    fonttitle=\bfseries,
    colback=purple!4,
    colframe=purple!50!black,
    boxrule=0.8pt,
    arc=2mm
}

\newtcolorbox{exercise}[1]{
    enhanced,
    breakable,
    title={Ejercicio: #1},
    fonttitle=\bfseries,
    colback=gray!5,
    colframe=gray!60!black,
    boxrule=0.8pt,
    arc=2mm
}

\newtcolorbox{note}{
    enhanced,
    breakable,
    title={Nota},
    fonttitle=\bfseries,
    colback=white,
    colframe=gray!50,
    boxrule=0.6pt,
    arc=2mm
}

% ==========================================================
% CONFIGURACIÓN FINAL
% ==========================================================
\renewcommand{\arraystretch}{1.25}

% ==========================================================
% DOCUMENTO
% ==========================================================
\begin{document}

% ---------------- PORTADA ----------------
\begin{titlepage}

\thispagestyle{empty}

\begin{center}

\vspace*{1cm}

{\large \textsc{\universidad}}

\vspace{0.2cm}

{\large \textsc{\facultad}}

\vspace{0.2cm}

{\large \carrera}

\vspace{3cm}

{\Huge \textbf{\materia}}

\vspace{1cm}

{\Large \titulo}

\vspace{0.8cm}

\textit{Comprender cada concepto es el primer paso para convertir el estudio en conocimiento.}

\vspace{1.2cm}

\rule{0.70\textwidth}{0.5pt}

\vspace{0.5cm}

{\large Apuntes de estudio}

\vspace{0.5cm}

\rule{0.70\textwidth}{0.5pt}

\vfill

{\large \autor}

\vspace{0.3cm}

{\large \anio}

\end{center}

\vfill

\begin{flushleft}
\textit{Este es un modesto aporte a los estudiantes de ``\materia'' no pretende sustituir el material oficial de la materia.}
\end{flushleft}

\end{titlepage}

% ---------------- ÍNDICE ----------------
\tableofcontents

% ==========================================================
% CAPÍTULO 1
% ==========================================================
\chapter[Organización del taller y cuestiones previas]{Organización del taller y cuestiones previas}

\section{Plan del taller}

El taller, que por las circunstancias de la clase se desarrolló de manera virtual en lugar de presencial, tiene por finalidad que los estudiantes realicen un \textbf{role playing} de una \textbf{audiencia preliminar}. La audiencia simulada comienza con la fijación de los hechos, continúa con la orden de la jueza sobre qué prueba se admite y culmina con la invitación a conciliar.

La clase se organizó en tres momentos:

\begin{enumerate}
    \item Consultas de los estudiantes acerca de los roles que cumplirán en la actuación.
    \item Exposición de la docente, con apoyo de una presentación, sobre la audiencia preliminar.
    \item Desarrollo del role playing de la audiencia preliminar.
\end{enumerate}

Durante la simulación se admite equivocarse: el ejercicio tiene una finalidad de aprendizaje y las correcciones se realizan al final.

\section{Roles de secretaría y mesa de entradas}

Cuando intervienen varias secretarias o secretarios, las funciones pueden distribuirse entre ellos:

\begin{itemize}
    \item Una secretaria \textbf{dialoga con la jueza} y, cuando esta lo requiere, aporta la agenda para fijar la audiencia de prueba testimonial.
    \item Otra secretaria \textbf{informa a las partes que la audiencia es videograbada}.
    \item La \textbf{mesa de entradas} solicita los documentos de identidad, identifica a los comparecientes y convoca la audiencia indicando los autos caratulados, donde siempre figura primero el apellido y luego el nombre, y el objeto del juicio (en este caso, desalojo por causales).
\end{itemize}

\section{Pliego de absolución de posiciones}

El pliego de posiciones debe presentarse, si la prueba confesional fue ofrecida, \textbf{hasta media hora antes de la audiencia}. Si se presenta con posterioridad a ese momento, el pliego «se murió». No corresponde, por lo tanto, presentar el pliego una vez que se frustra la conciliación.

\begin{note}
En la práctica descripta en clase, el Código establece la presentación del pliego por escrito porque el sistema era en formato papel, pero las posiciones suelen formularse \textbf{a viva voz} en la propia audiencia. Además, la absolución de posiciones se practica poco: los abogados de ambas partes suelen acordar no tomar posiciones y los clientes llegan preparados, de modo que esta prueba \textbf{cayó en desuso}.
\end{note}

% ==========================================================
% CAPÍTULO 2
% ==========================================================
\chapter[La audiencia preliminar como «caballo de Troya»]{La audiencia preliminar como «caballo de Troya» del proceso}

\section{La metáfora del caballo de Troya}

La docente relaciona la audiencia preliminar con la historia del caballo de Troya, que presenta como una \textbf{estrategia}: los soldados griegos se ocultaron dentro de un caballo de madera, lo llevaron a la puerta de la ciudad, los troyanos lo ingresaron a la plaza central y, de noche, los soldados salieron y tomaron la ciudad, ganando la guerra.

La audiencia preliminar es igualmente algo \textbf{estratégico} dentro del proceso: es un momento culminante en la fijación del conflicto, de los hechos y de los medios de prueba que serán admitidos.

\section{Inmediación y plazo razonable}

\begin{definition}{Inmediación en la audiencia preliminar}
Es la primera y única posibilidad del proceso en que la jueza o el juez se encuentra con los abogados, las abogadas y las partes, de manera presencial o virtual, de modo que se produce un intercambio de ida y vuelta entre todos los intervinientes procesales.
\end{definition}

En este encuentro el proceso comienza a revelar qué puede solucionarse del conflicto. El conflicto puede solucionarse o, llegado el momento final, ser resuelto por el juez; la docente recuerda al respecto su esquema del desalojo, que culmina con una balanza de la justicia compuesta por personas.

La audiencia se plantea como estrategia para hacer frente a la condena por los \textbf{plazos no razonables} del proceso civil. El proceso civil suele excederse en el tiempo y, cuando la justicia llega, es una justicia tardía. Argentina fue condenada por la Corte Interamericana de Derechos Humanos por no cumplir con el plazo razonable. Las herramientas creadas por el legislador permiten solucionar la conflictiva o resolverla de forma más acotada y depurada.

\begin{important}
La audiencia preliminar es una estrategia: es la primera y única ocasión de inmediación del proceso, en la cual el conflicto puede solucionarse o resolverse de modo más acotado y depurado, para cumplir con el plazo razonable que el proceso civil suele exceder.
\end{important}

\section{La audiencia en la práctica: delegación y justicia tardía}

En la práctica, muy pocos juzgados toman la audiencia como un «caballo de Troya» del proceso. Con frecuencia la audiencia es tomada por un empleado o «audiencista» del juzgado, que la despacha a la ligera preguntando si hay conciliación y, si no la hay, pasa a la etapa siguiente («ya proveerán»). Esta no es la función para la cual el legislador estableció la audiencia.

El Código establece el \textbf{carácter indelegable} de la audiencia; sin embargo, la generalidad de los casos no lo cumple, no existe sanción para el magistrado y los abogados no suelen presentar denuncias ante el Consejo de la Magistratura.

La docente recuerda que en las jornadas se mencionó que ingresan 140.000 causas por año, dato que cita como contexto para señalar que se desperdicia esta oportunidad.

\section{El momento único de inmediación: la jueza como directora del proceso}

El momento del artículo 360 es el único en que la jueza conoce, presencial, física o virtualmente, el rostro de los sujetos en conflicto bajo su jurisdicción, así como el de los abogados y abogadas, aunque ya los conozca de otros procesos, para ver cómo trabajan.

La jueza es \textbf{directora del proceso}: debe procurar que las partes solucionen el conflicto o acotar los plazos para cumplir con la Corte Interamericana de Derechos Humanos y resolver en un tiempo razonable y conforme a derecho.

\section{Roles de la jueza, de los abogados y de las partes}

Un error frecuente de los abogados es indicar al cliente que no hable en la audiencia, o que no concurra y otorgue un poder para que concurra solo el apoderado. La parte es quien estuvo en el conflicto y quien resultó lesionada; la jueza no puede trabajar bien sin tenerla presente, de manera virtual o presencial. No es lo mismo la perspectiva del abogado, externo al conflicto, que la de la parte.

La participación de las partes es importante porque pueden reflotar otros temas, lo que conduce al camino de la conciliación.

\subsection{El valor de la palabra}

Para ilustrar este punto la docente recurre a la imagen de una narradora de cuentos, abogada de estudio, que recorre el país llevando libros mediante donaciones. Esta narradora sostiene que la palabra es maravillosamente sanadora porque pasa por el corazón y genera comunicación directa entre las personas.

Aplicado a la audiencia: no es lo mismo que la jueza trabaje con la parte, en persona o virtualmente, que trabaje solo con el abogado. La palabra debe tener apertura, siempre con respeto. Es más provechosa la presencia de las partes en el despacho, pero si no es posible se hace de manera virtual. El resultado es más rico si las partes concurren y los abogados les permiten expresarse con respeto.

\section{Función polifuncional de la audiencia: el artículo 360}

La audiencia cumple una función \textbf{polifuncional}: no es solo conciliar y no es solo oralidad. Según los incisos del artículo 360, en ella:

\begin{itemize}
    \item se invita a las partes a una conciliación;
    \item se fijan los hechos controvertidos y conducentes;
    \item se fijan los medios probatorios admitidos y ofrecidos por las partes;
    \item se recibe la prueba confesional, si corresponde;
    \item se fija la prueba testimonial, si fue ofrecida y admitida;
    \item se sortean los peritos, si fueron ofrecidos y admitida la prueba;
    \item se escucha a las partes que se opongan a la apertura a prueba;
    \item la jueza puede declarar la causa de puro derecho (inciso 6), caso en el cual no hay apertura a medios de prueba y solo resta determinar qué derecho se aplica a la actora y a la demandada; el expediente pasa entonces para sentencia definitiva.
\end{itemize}

\begin{formula}{Artículo 360 CPCCN --- Audiencia preliminar}
A los fines del artículo precedente el juez citará a las partes a una audiencia, que presidirá, con carácter indelegable. Si el juez no se hallare presente no se realizará la audiencia, debiéndose dejar constancia en el libro de asistencia.

1) Invitará a las partes a una conciliación o a encontrar otra forma de solución de conflictos que acordarán en la audiencia. El juez podrá, si la naturaleza y el estado del conflicto lo justifican, derivar a las partes a mediación. En este supuesto, se suspenderá el procedimiento por treinta (30) días contados a partir de la notificación del mediador a impulso de cualquiera de las partes. Vencido este plazo, se reanudará el procedimiento a pedido de cualquiera de las partes, lo que dispondrá el juez sin sustanciación, mediante auto que se notificará a la contraria.

2) Recibirá las manifestaciones de las partes con referencia a lo prescripto en el artículo 361 del presente Código, debiendo resolver en el mismo acto.

3) Oídas las partes, fijará los hechos articulados que sean conducentes a la decisión del juicio sobre los cuales versará la prueba.

4) Recibirá la prueba confesional si ésta hubiera sido ofrecida por las partes. La ausencia de uno o de todos los absolventes no impedirá la celebración de la audiencia preliminar.

5) Proveerá en dicha audiencia las pruebas que considere admisibles y concentrará en una sola audiencia la prueba testimonial, la que se celebrará con presencia del juez en las condiciones establecidas en este capítulo. Esta obligación únicamente podrá delegarse en el secretario o, en su caso, en el prosecretario letrado.

6) Si correspondiere, decidirá en el acto de la audiencia que la cuestión debe ser resuelta como de puro derecho con lo que la causa quedará concluida para definitiva.
\end{formula}

% TODO: la docente enumeró los incisos de memoria; el texto del artículo figura en las notas originales con la indicación de cotejarlo con el texto oficial vigente.
\begin{note}
Las notas originales indican que la docente enumeró los incisos de memoria durante la exposición y que conviene cotejar la redacción transcripta con el texto oficial vigente antes de citarla.
\end{note}

\section{Audiencia activa o audiencia de cinco minutos}

La audiencia es un momento «riquísimo» y muy activo del proceso si se realiza activamente; en cambio, puede durar cinco minutos cuando la toma cualquier miembro del tribunal y se limita a preguntar si hubo conciliación y, ante la respuesta negativa, levantar el acta y dejar el despacho de la prueba para después. De esa manera se pierde todo el debate y la riqueza de la inmediación.

\begin{table}[H]
\centering
\caption{La audiencia preliminar según el Código y en la práctica judicial}
\small
\begin{tabularx}{\textwidth}{>{\raggedright\arraybackslash}p{2.6cm} >{\raggedright\arraybackslash}X >{\raggedright\arraybackslash}X}
\toprule
\textbf{Aspecto} & \textbf{Lo que establece el Código} & \textbf{Lo que la docente describe de la práctica} \\
\midrule
Quién la toma & La jueza o el juez, con carácter indelegable (art. 360). & A veces la toma un empleado o «audiencista» del juzgado. \\
Conciliación & Se invita a las partes a conciliar y se trabaja el conflicto. & Se pregunta «¿hay conciliación?»; si no la hay, se pasa a la etapa siguiente. \\
Hechos y prueba & Se fijan los hechos conducentes y se depuran los medios con las partes presentes. & Se deja todo para después: «ya proveerán». \\
Duración & Momento «riquísimo» y muy activo del proceso. & Puede durar cinco minutos. \\
Sanción & El Código dispone el carácter indelegable. & No hay sanción para el magistrado y los abogados no suelen denunciar ante el Consejo de la Magistratura. \\
\bottomrule
\end{tabularx}
\end{table}

\begin{note}
Ante la pregunta de un estudiante sobre si, de no realizarse la audiencia del artículo 360, se acepta toda la prueba ofrecida, la docente responde que no: la prueba debe fijarse en esa audiencia.
\end{note}

% ---------------- CORNELL CAPÍTULO 2 ----------------
\section{Cuadro Cornell: la audiencia preliminar}

\begingroup
\small
\begin{longtable}{
    >{\raggedright\arraybackslash}p{0.26\textwidth}
    >{\raggedright\arraybackslash}p{0.66\textwidth}
}
\toprule
\textbf{Preguntas / palabras clave} & \textbf{Notas / respuestas} \\
\midrule
\endfirsthead
\toprule
\textbf{Preguntas / palabras clave} & \textbf{Notas / respuestas} \\
\midrule
\endhead
\bottomrule
\endfoot
\textbf{¿Por qué se compara la audiencia preliminar con el «caballo de Troya»?} &
Porque, como los soldados griegos ocultos en el caballo, la audiencia es una estrategia: dentro de ella está la posibilidad de transformar la estructura del conflicto, solucionarlo o resolverlo de manera más acotada y depurada, en un plazo razonable. \\[0.6em]

\textbf{¿Qué es la inmediación y por qué se la llama «única»?} &
Es el momento en que la jueza o el juez se encuentra, de forma presencial o virtual, con los abogados y las partes. Es la primera y única vez del proceso en que hay una comunicación de ida y vuelta entre todos los intervinientes. \\[0.6em]

\textbf{¿Qué problema busca atenuar la audiencia?} &
El plazo no razonable del proceso civil, por el que Argentina fue condenada por la Corte Interamericana de Derechos Humanos. Usada con seriedad, evita una «justicia tardía». \\[0.6em]

\textbf{¿Qué papel cumplen la jueza, los abogados y las partes?} &
La jueza dirige el proceso y procura que las partes solucionen el conflicto o se acorten los plazos. Los abogados asisten a las partes sin silenciarlas. Las partes deben estar presentes: su palabra puede reabrir temas y conducir a la conciliación. \\[0.6em]

\textbf{¿Qué se hace en la audiencia del artículo 360?} &
Es una audiencia polifuncional: invita a conciliar, resuelve oposiciones a la apertura a prueba, fija los hechos conducentes y controvertidos, admite los medios de prueba, recibe la confesional si fue ofrecida, fija la testimonial, sortea peritos y puede declarar la causa de puro derecho (inciso 6). \\[0.6em]

\textbf{¿Puede delegarse la audiencia?} &
El Código la declara indelegable, pero en la práctica a veces la toma un empleado y dura minutos. No hay sanción para el magistrado. \\
\midrule
\multicolumn{2}{p{0.92\textwidth}}{\textbf{Resumen:} la audiencia preliminar del artículo 360 es el único momento del proceso en que el juez conoce a las partes y a sus abogados. Usada como estrategia, permite solucionar el conflicto o resolverlo de forma acotada y depurada, atenuando el problema de los plazos no razonables. Por eso es indelegable y polifuncional, aunque en la práctica suele despacharse a la ligera.} \\
\end{longtable}
\endgroup

% ==========================================================
% CAPÍTULO 3
% ==========================================================
\chapter[Prueba, impugnación y recursos]{Prueba, impugnación y recursos}

\section{La resolución que fija los hechos y admite la prueba}

Después de la audiencia se dicta una resolución, visible en el expediente electrónico, que establece \textbf{cuáles hechos son conducentes y controvertidos} y qué \textbf{medios de prueba fueron admitidos}.

\section{Artículo 379: qué se puede impugnar y qué no}

Las partes y sus abogados solo pueden impugnar cuando la jueza \textbf{rechaza un hecho articulado} por las partes en la demanda, la contestación o la reconvención para que sea objeto de prueba. No pueden impugnar, en cambio, cuando la jueza \textbf{no admite un medio de prueba} (por ejemplo, una pericial médica oncológica): por el artículo 379 esa resolución es inapelable. Ese medio de prueba debe replantearse en la etapa del artículo 260 del Código Procesal, es decir, si se apela la sentencia definitiva y la apelación se concede con efecto libre.

El rechazo de un hecho afirmado se impugna a viva voz en la propia audiencia (apelación \textit{in voce}).

\begin{important}
Hechos sí, medios de prueba no. Las partes pueden impugnar el rechazo de un hecho articulado en la demanda, la contestación o la reconvención; no pueden apelar el rechazo de un medio de prueba (art. 379), que solo se replantea ante la cámara según el art. 260.
\end{important}

\begin{formula}{Artículo 379 CPCCN --- Inapelabilidad de las resoluciones sobre prueba}
Serán inapelables las resoluciones del juez sobre producción, denegación y sustanciación de las pruebas; y si se hubiere negado alguna medida probatoria, la parte interesada podrá solicitar a la cámara que la diligencie cuando el expediente le fuera remitido para conocer del recurso concedido contra la sentencia definitiva.
\end{formula}

\begin{formula}{Artículo 260, inciso 2 CPCCN --- Replanteo de la prueba en segunda instancia}
Dentro del quinto día de notificada la providencia a que se refiere el artículo anterior y en un solo escrito, las partes deberán: [\ldots] 2) Indicar las medidas probatorias denegadas en primera instancia o respecto de las cuales hubiese mediado declaración de negligencia, que tengan interés en replantear en los términos de los artículos 379 y 385 in fine. La petición será fundada, y resuelta sin sustanciación alguna.
\end{formula}

\begin{note}
La docente remite al artículo 260 sin leerlo; el inciso pertinente figura transcripto en las notas originales. El tramo final del artículo 379 se completó conforme a la lectura realizada en la clase (ver el capítulo sobre el simulacro); conviene cotejar ambos textos con la redacción oficial.
\end{note}

\begin{table}[H]
\centering
\caption{Qué se puede impugnar y qué no (art. 379)}
\small
\begin{tabularx}{\textwidth}{>{\raggedright\arraybackslash}p{4.4cm} >{\raggedright\arraybackslash}p{2.6cm} >{\raggedright\arraybackslash}X}
\toprule
\textbf{Situación} & \textbf{¿Se puede impugnar?} & \textbf{Qué hacer} \\
\midrule
La jueza rechaza un hecho articulado en la demanda, la contestación o la reconvención & Sí & Impugnar a viva voz en la propia audiencia. \\
La jueza no admite un medio de prueba (por ejemplo, una pericial médica oncológica) & No (art. 379) & Replantear ante la cámara (art. 260) si se apela la sentencia definitiva y se concede libremente. \\
Proceso sumarísimo & Apelación muy limitada (art. 498) & No procede replantear la prueba por el art. 379: «estar en el horno». \\
\bottomrule
\end{tabularx}
\end{table}

\section{Proceso sumarísimo y recursos}

La etapa recursiva no se alcanza a repasar en la clase y debe estudiarse por separado. Cuando se apela una sentencia dictada en un juicio de procedimiento ordinario, la prueba denegada puede replantearse en los términos del artículo 379. En el \textbf{proceso sumarísimo}, en cambio, la prueba perdida no puede replantearse por el artículo 379.

El desalojo tramita a menudo por proceso sumarísimo, cuya apelación es mucho más limitada y cuyas excepciones oponibles son pocas. La docente indica leer con atención el artículo 498 y señala que este tema suele ser objeto de preguntas en el examen final.

\begin{formula}{Artículo 498, inciso 6 CPCCN --- Proceso sumarísimo: apelación}
Sólo serán apelables la sentencia definitiva y las providencias que decreten o denieguen medidas precautorias. La apelación se concederá en relación, en efecto devolutivo, salvo [\ldots]
\end{formula}

% TODO: contenido incompleto o ambiguo en las notas originales (la transcripción del artículo 498 se interrumpe en «salvo [...]»)
\begin{note}
Las notas originales transcriben solo el primer tramo del inciso. La docente remite al artículo 498 («léanselo bien») porque el desalojo tramita muy a menudo por proceso sumarísimo.
\end{note}

La docente ilustra la importancia práctica con un caso: si, al recibirse, se interpone en un primer desalojo el recurso concedido libremente y se solicita el replanteo del artículo 379 en un proceso sumarísimo, la prueba se pierde.

\section{El esquema del proceso y la homologación del acuerdo}

El esquema del proceso, que suele ponerse en el pizarrón, ubica las normas del Código que se desarrollan en la audiencia preliminar junto con otras normas del Código Procesal.

Si en la audiencia las partes llegan a un acuerdo, la jueza labra un \textbf{acta} y debe \textbf{homologar} el acuerdo mediante una resolución, que se incorpora al expediente electrónico. Aunque las partes lleguen a un acuerdo y la jueza esté presente, la jueza igualmente debe dictar la resolución que homologa el acuerdo alcanzado en la audiencia del artículo 360.

\section{Los dos momentos para declarar la causa de puro derecho}

\begin{definition}{Causa de puro derecho}
Es aquella en la cual el juez o jueza debe determinar el derecho aplicable a los hechos, sin que se discutan medios de prueba, porque no hay prueba por producir: la prueba ya fue producida como prueba anticipada, las partes desistieron de los medios o solo existe prueba documental.
\end{definition}

La causa puede declararse de puro derecho en dos momentos:

\begin{enumerate}
    \item en el \textbf{artículo 359}, una vez trabada la litis desde el punto de vista sustancial;
    \item en la propia audiencia, según el \textbf{artículo 360, inciso 6}.
\end{enumerate}

\begin{formula}{Artículo 359 CPCCN --- Puro derecho o apertura a prueba}
Contestado el traslado de la demanda o reconvención, en su caso, o vencidos los plazos para hacerlo, resueltas las excepciones previas, si la cuestión pudiera ser resuelta como de puro derecho, así se decidirá y, firme que se encuentre la providencia, se llamará autos para sentencia. Si se hubiesen alegado hechos conducentes acerca de los cuales no hubiese conformidad entre las partes, aunque éstas no lo pidan, el juez recibirá la causa a prueba procediendo de acuerdo con lo preceptuado en el artículo 360.
\end{formula}

\begin{table}[H]
\centering
\caption{Dos momentos para declarar la causa de puro derecho}
\small
\begin{tabularx}{\textwidth}{>{\raggedright\arraybackslash}p{4.4cm} >{\raggedright\arraybackslash}p{2.6cm} >{\raggedright\arraybackslash}X}
\toprule
\textbf{Momento} & \textbf{Norma} & \textbf{Qué ocurre} \\
\midrule
Una vez trabada la litis desde el punto de vista sustancial & Art. 359 & Se decide que la cuestión es de puro derecho antes de la audiencia. \\
Dentro de la propia audiencia preliminar & Art. 360, inc. 6 & La jueza lo declara en el acto; la causa queda concluida para definitiva. \\
Efecto común & Ambos & Ya no se discuten medios de prueba, sino el derecho aplicable a los hechos. \\
\bottomrule
\end{tabularx}
\end{table}

\section{Fijación de la testimonial y constancia del acuerdo}

En la audiencia la jueza fija la prueba testimonial: señala las fechas de audiencia para los testigos de la actora y de la demandada. Luego otorga un espacio de reflexión a las partes, que pueden terminar en un acuerdo.

El acuerdo debe ser \textbf{descripto}: la secretaria de la jueza toma nota de su contenido. Como la audiencia está videograbada, el acuerdo queda plasmado a viva voz en las cámaras y la grabación se sube luego al expediente. Lo que se firma es un \textbf{acta de comparecencia}; a veces ni siquiera se firma, y la comparecencia queda registrada en la videograbación.

% ---------------- CORNELL CAPÍTULO 3 ----------------
\section{Cuadro Cornell: prueba, recursos y acuerdo}

\begingroup
\small
\begin{longtable}{
    >{\raggedright\arraybackslash}p{0.26\textwidth}
    >{\raggedright\arraybackslash}p{0.66\textwidth}
}
\toprule
\textbf{Preguntas / palabras clave} & \textbf{Notas / respuestas} \\
\midrule
\endfirsthead
\toprule
\textbf{Preguntas / palabras clave} & \textbf{Notas / respuestas} \\
\midrule
\endhead
\bottomrule
\endfoot
\textbf{¿Qué se puede impugnar y qué no? (art. 379)} &
Se puede impugnar el rechazo de un hecho articulado en la demanda, la contestación o la reconvención. No se puede apelar el rechazo de un medio de prueba: se replantea ante la cámara (art. 260) si se apela la sentencia definitiva y se concede libremente. \\[0.6em]

\textbf{¿Qué ocurre en el proceso sumarísimo? (art. 498)} &
El desalojo tramita muy a menudo por proceso sumarísimo, con una apelación mucho más limitada. No hay replanteo de la prueba por el art. 379. Es tema frecuente de las preguntas del examen final. \\[0.6em]

\textbf{¿Cuándo se declara la causa de puro derecho?} &
En dos momentos: el art. 359 (trabada la litis) y el art. 360, inciso 6 (dentro de la audiencia). El juez define el derecho aplicable porque no hay medios de prueba por producir. \\[0.6em]

\textbf{¿Cómo se concreta un acuerdo?} &
Se describe en la audiencia y queda registrado en el acta o en la videograbación subida al expediente. La jueza debe homologarlo mediante una resolución que impacta en el expediente electrónico. \\[0.6em]

\textbf{¿Hasta cuándo se presenta el pliego de posiciones?} &
Hasta media hora antes de la audiencia; después «se murió». En la práctica las posiciones se formulan a viva voz y la absolución cayó en desuso. \\
\midrule
\multicolumn{2}{p{0.92\textwidth}}{\textbf{Resumen:} en materia de prueba, las partes pueden impugnar el rechazo de un hecho pero no el de un medio de prueba (art. 379), que solo se replantea ante la cámara (art. 260); en el desalojo, que suele ser sumarísimo, la apelación es todavía más limitada (art. 498). La causa puede declararse de puro derecho en el art. 359 o en el art. 360, inciso 6. El acuerdo alcanzado en la audiencia queda registrado en acta o videograbación y debe ser homologado por la jueza.} \\
\end{longtable}
\endgroup

% ==========================================================
% CAPÍTULO 4
% ==========================================================
\chapter[Simulacro de la audiencia preliminar]{Simulacro de la audiencia preliminar: «Riquel c/ Milán s/ desalojo por cambio de destino»}

\section{Instrucciones del role playing}

La jueza dirige la audiencia. Las secretarias son sus colaboradoras: si la jueza omite un paso o avanza demasiado rápido, deben intervenir con naturalidad, como si estuvieran en la audiencia real (por ejemplo, indicando «su señoría, vamos a\ldots»). Las docentes permanecen en silencio aunque se cometan errores; no se interrumpe el desarrollo y las observaciones se realizan al final.

\begin{note}
La docente menciona un canal de YouTube con audiencias simuladas grabadas, como recurso audiovisual complementario para estudiar el desarrollo de la audiencia.
\end{note}

\section{Apertura, identificación de las partes y constancias formales}

La secretaría abre la audiencia indicando la hora y la fecha, el expediente y los autos caratulados (\textit{Riquel Juan Román y otros contra Milán Martín y otros sobre desalojo por cambio de destino}), tramitados ante un juzgado nacional de primera instancia en lo civil. Antes de comenzar se verifica la identidad y comparecencia de las partes, de sus letrados y de los demás intervinientes, comenzando por la parte actora y siguiendo con la parte demandada y el fiador.

Cada compareciente exhibe su documento de identidad ante la cámara; los letrados indican nombre, apellido, matrícula profesional y el carácter en el que intervienen (por ejemplo, abogados patrocinantes). Como el fiador no se encontraba presente, se solicitó a su letrado que se identificara.

Se deja constancia de la comparecencia de las partes, de sus letrados y del fiador individualizado, y se informa a los presentes que la audiencia es íntegramente grabada y que su registro se incorporará al expediente judicial electrónico. Cumplidas las identificaciones y constancias, la secretaría da intervención a la jueza, que dirige la audiencia.

\section{Apertura de la jueza y exposición de la parte actora}

La jueza expone que preside la audiencia del artículo 360 con carácter indelegable y fija el orden de intervención: primero la parte actora, luego la parte demandada, después el fiador y por último los letrados patrocinantes de cada parte.

La parte actora relata los hechos:

\begin{itemize}
    \item El inmueble fue heredado y se encuentra alquilado al demandado desde hace un año, por un canon de 750.000 pesos, para que viva allí con su familia.
    \item Se enteraron, por mensajes de vecinos y videos en redes sociales, de que se lo utiliza para otros fines: hay olores, música, movimiento de colectivos y de gente, ruidos y suciedad.
    \item Sostienen que el destino pactado en el contrato era habitacional, que el emprendimiento daña el inmueble (en particular la cocina) y que no obtuvieron respuesta cuando intentaron comunicarse.
    \item Solicitan la restitución del inmueble.
\end{itemize}

\section{Exposición de la parte demandada y de los letrados}

La parte demandada reconoce que tiene un emprendimiento pequeño de venta de repostería, iniciado tras perder el empleo para sostener a la familia. Explica que la venta es normalmente en línea y que, excepcionalmente, con motivo de un partido importante, tuvo una gran demanda y debió ampliar la venta a trato directo. Sostiene que fue un caso puntual y temporal, que intentó comunicarse con los locadores por varios medios sin obtener respuesta y que se tramita un permiso para vender en un local aparte, fuera del domicilio.

Los letrados de la parte actora agregan que:

\begin{itemize}
    \item el titular del derecho de habitar no vive en el inmueble, que solo ocupa la dueña del emprendimiento;
    \item el alquiler fue para uso habitacional y el uso comercial está prohibido;
    \item los hechos se acreditan con videos, mensajes, conversaciones y fotos incorporados a la demanda;
    \item ante la imposibilidad de comunicarse con el demandado, lo intimaron y fue debidamente notificado, lo cual consta en el expediente.
\end{itemize}

Los letrados de la parte demandada sostienen que:

\begin{itemize}
    \item la situación del emprendimiento es excepcional y se prioriza la venta en línea; en el domicilio solo hay un depósito de las mercaderías, sin venta directa;
    \item la concurrencia de personas fue algo de una sola vez;
    \item se trata de una cuestión temporal debida a un trámite administrativo del gobierno de la ciudad, que no habría otorgado aún el permiso para vender en un lugar habilitado;
    \item los alquileres se encuentran pagos y lo conveniente es que las costas se impongan por su orden;
    \item existe prueba testimonial de otros vecinos.
\end{itemize}

\section{Fijación de los hechos conducentes y controvertidos}

Escuchadas las partes, la jueza fija los hechos controvertidos y conducentes de acuerdo con el \textbf{inciso 3 del artículo 360} y con los escritos postulatorios: el uso del inmueble y la verificación del cambio de destino alegado por la parte actora, esto es, que el inmueble fue alquilado únicamente con destino de vivienda y que esa cláusula no se estaría cumpliendo.

La letrada de la parte actora agrega que también existiría violación al orden público y a las normas de los vecinos. La jueza responde que fija los hechos conducentes conforme a los escritos postulatorios de las partes.

\section{Depuración y admisión de la prueba; oposición y artículo 379}

La jueza procede a la depuración de la prueba ofrecida:

\begin{itemize}
    \item \textbf{admite} la prueba documental (contrato de locación, cartas, documentos, fotografías, videos y capturas de chat) y la testimonial;
    \item \textbf{rechaza} la prueba de reconocimiento judicial y la pericial ofrecidas por las partes.
\end{itemize}

Si algún abogado se opone al rechazo, la jueza debe leer el \textbf{artículo 379} y explicar que, conforme a él, la resolución es inapelable. Luego se provee la prueba restante y se fija la audiencia de prueba testimonial para el 10 de diciembre, con la colaboración de la secretaría y su agenda.

\begin{note}
El texto completo del artículo 379 figura en el capítulo anterior. En el simulacro, ante la oposición de los abogados al rechazo de la prueba pericial, la jueza lee el artículo y explica que la resolución es inapelable.
\end{note}

\section{Invitación a conciliar y negociación entre las partes}

La jueza invita a las partes a resolver el conflicto por sí mismas. Señala que un conflicto no necesariamente debe ser resuelto por un tercero (el sistema judicial), que un acuerdo es más rápido, menos costoso y se adecúa más a los intereses de las partes, y que evita un trámite largo y costoso.

\textbf{Posición de la parte actora:}

\begin{itemize}
    \item no pretende cesar la relación contractual, pero solicita el cese total del cambio de destino;
    \item solicita el arreglo total de la cocina, con fotos del estado del inmueble por dentro;
    \item solicita que el contratante viva en el inmueble;
    \item no acepta prórroga ni extensión, y plantea como alternativas el cese del destino (cierre del emprendimiento) o la restitución del inmueble.
\end{itemize}

\textbf{Posición de la parte demandada:}

\begin{itemize}
    \item los muebles de la cocina pertenecen a los inquilinos y el inmueble se devolverá en las condiciones en que fue entregado;
    \item reconoce que se cambió en parte el destino por una cuestión económica derivada de la pérdida de trabajo, pero sostiene la importancia de respetar el contrato hasta su finalización;
    \item solicita un plazo razonable para tramitar el permiso y vender en un local adecuado, pues sin el emprendimiento no podrían pagar el canon locativo;
    \item propone un convenio aparte para dar ingresos a un familiar de la parte actora cuando se abra el local, y un reajuste del alquiler por un tiempo corto (tres o seis meses);
    \item ofrece reemplazar las ventas presenciales por envíos, e instalar extractores para evitar olores.
\end{itemize}

La parte actora cuestiona la instalación del extractor por temor a que se agujeree la pared; la parte demandada responde que solo requiere unos tornillos y que la salida de aire de la casa permite conectarlo sin modificar la estructura.

\section{Acuerdo, costas y cierre de la audiencia}

\begin{example}{Acuerdo alcanzado en el simulacro}
Las partes acordaron que se cambiará el lugar donde se fabrican las tortas; que, por un lapso de tres semanas, se permitirá que las personas sigan yendo a retirar la pastelería, pero no la producción; y que, respecto de los ruidos, se colocará doble vidrio para aislar el sonido que emitan tanto los niños como los clientes.

Respecto de las \textbf{costas}, la parte actora propuso inicialmente que los honorarios los pagara la demandada; la parte demandada propuso que cada parte abone las suyas («por su orden»), por las condiciones en que se dieron los hechos, y todos estuvieron de acuerdo.
\end{example}

La jueza deja constancia de que el acuerdo constará en un acta, pasará a homologación y será subido al sistema del Poder Judicial de la Nación.

% TODO: contenido incompleto o ambiguo en las notas originales (la parte demandada cuestionó por qué estaría prohibida la producción de tortas en su domicilio; el apunte no registra respuesta clara)

% ==========================================================
% CAPÍTULO 5
% ==========================================================
\chapter[Devolución sobre el simulacro]{Devolución de las docentes sobre el simulacro}

\section{Fórmulas conciliatorias propuestas por el juez}

Un aspecto sustancial es que, en la audiencia, \textbf{el juez o la jueza puede proponer fórmulas conciliatorias}. En el simulacro, dado el intercambio de las partes, no parecía que se fuera a llegar a un acuerdo, y la jueza podría haber intervenido.

El acuerdo conciliatorio al que se arriba en un juzgado es una forma normal de terminación del proceso. A diferencia de la mediación privada, el juez puede dirigir la conciliación proponiendo soluciones concretas (por ejemplo, que el demandado permanezca un mes más y luego cese la actividad) y preguntando a la contraparte si le parecen prudentes. También puede señalar qué considerará al dictar sentencia, a la manera de un conciliador.

En cuanto a la prueba:

\begin{itemize}
    \item la testimonial se fija en otro momento;
    \item se realiza una depuración de la prueba, si se hace lugar a ella y es conducente;
    \item la confesional no fue ofrecida en el simulacro; de haberlo sido, podría haberse invitado a las partes a reconocer si existían o no los hechos.
\end{itemize}

% TODO: contenido incompleto o ambiguo en las notas originales (observación sobre la identificación de los abogados por «tomo y folio» en Capital)
\begin{note}
La docente observa que, en la Capital, los abogados no se identifican por una matrícula sino por \textit{tomo y folio} (por ejemplo, «tomo dos, folio 540»), y agrega que no es relevante.
\end{note}

\section{Integración de la litis y rol del fiador}

\textbf{La litis no estaba integrada}: una de las personas demandadas no concurrió. Además, el fiador es parte y no se le permitió hablar. El fiador podría haber propuesto, por ejemplo, un acuerdo para que dentro de dos meses la parte demandada se retire por el cambio de destino, ya que se compromete como fiador del cumplimiento de las obligaciones del contrato. Si no concurren todas las partes, es difícil llegar a un arreglo, porque deben estar todas presentes y la litis debe integrarse.

\begin{important}
Sin litis integrada es difícil acordar. Para llegar a un arreglo deben estar todas las partes; si el fiador estuvo notificado y no concurrió, se le notifica el acuerdo para que preste conformidad u oposición y se resguarde su defensa en juicio.
\end{important}

\section{Si el fiador no concurre en la práctica judicial}

En la realidad judicial, si el fiador estuvo debidamente notificado y no concurrió, la audiencia igualmente se celebra y el acuerdo se celebra, pero la jueza dispone previa notificación al fiador, si está presentado en el proceso, para que preste oposición o conformidad. El acuerdo se notifica al fiador y queda cerrado una vez vencido el plazo. Esto se funda en el respeto de los principios procesales de la defensa en juicio y de la igualdad de las partes: el fiador es, en definitiva, el más perjudicado si el inquilino incumple y se provoca un desalojo.

\section{Conciliación judicial y mediación}

El juez procura la conciliación para terminar el proceso por una forma normal y tener un \textbf{expediente menos}; por ello se le otorga ese poder. En ese marco la docente cita nuevamente los 140.000 ingresos anuales de causas y señala que una jueza debería dedicarse a la audiencia preliminar «para limpiar», porque de lo contrario el volumen de causas es interminable. En la práctica, tanto en sede laboral como civil, el juez insiste en que se llegue a un acuerdo.

\begin{table}[H]
\centering
\caption{Acuerdo judicial y mediación privada (según lo expuesto en clase)}
\small
\begin{tabularx}{\textwidth}{>{\raggedright\arraybackslash}p{3cm} >{\raggedright\arraybackslash}X >{\raggedright\arraybackslash}X}
\toprule
\textbf{Aspecto} & \textbf{Acuerdo conciliatorio en el juzgado} & \textbf{Mediación privada} \\
\midrule
Quién conduce & El juez puede proponer fórmulas y dirigir como conciliador. & El mediador. \\
Homologación & El juez debe homologar el acuerdo. & Según lo dicho, el acuerdo del mediador no necesita homologación. \\
Interés del juez & Terminar el proceso por una forma normal y tener un expediente menos. & --- \\
\bottomrule
\end{tabularx}
\end{table}

\section{Observaciones sobre el caso simulado}

En la audiencia simulada no había hechos controvertidos, porque ninguna parte negó lo afirmado por la otra; el desacuerdo se refería al cierre del emprendimiento y a un posible plazo. Si no hubiera habido acuerdo, el caso podría haber seguido el trámite directo o haberse declarado de puro derecho, sin necesidad de convocar a la audiencia preliminar. La docente agrega que las pruebas documentales no estaban certificadas por escribano, lo que complicaba el planteo.

% TODO: contenido incompleto o ambiguo en las notas originales (diálogo entre las docentes sobre trámite directo, puro derecho y documental no certificada)

% ---------------- CORNELL CAPÍTULOS 4 Y 5 ----------------
\section{Cuadro Cornell: simulacro y devolución}

\begingroup
\small
\begin{longtable}{
    >{\raggedright\arraybackslash}p{0.26\textwidth}
    >{\raggedright\arraybackslash}p{0.66\textwidth}
}
\toprule
\textbf{Preguntas / palabras clave} & \textbf{Notas / respuestas} \\
\midrule
\endfirsthead
\toprule
\textbf{Preguntas / palabras clave} & \textbf{Notas / respuestas} \\
\midrule
\endhead
\bottomrule
\endfoot
\textbf{¿Qué hechos fijó la jueza en el simulacro?} &
Conforme al art. 360, inciso 3 y a los escritos postulatorios: el uso del inmueble y la verificación del cambio de destino alegado por la actora (cláusula de destino exclusivo de vivienda). \\[0.6em]

\textbf{¿Qué prueba admitió y cuál rechazó?} &
Admitió la documental (contrato, cartas, documentos, fotografías, videos y capturas de chat) y la testimonial. Rechazó el reconocimiento judicial y la pericial. Ante la oposición, leyó el art. 379: la resolución es inapelable. \\[0.6em]

\textbf{¿Qué puede hacer el juez para conciliar?} &
Proponer fórmulas y dirigir como conciliador, señalando qué tendrá en cuenta al sentenciar. El acuerdo judicial requiere homologación; el juez busca terminar el proceso y tener un expediente menos. \\[0.6em]

\textbf{¿Por qué importa el fiador?} &
Es parte: la litis debe estar integrada. Si estaba notificado y no concurrió, se le notifica el acuerdo para que preste conformidad u oposición y se resguarde la defensa en juicio. \\[0.6em]

\textbf{¿Cómo se diferencia el acuerdo judicial de la mediación privada?} &
En el juzgado el juez puede proponer fórmulas y dirigir como conciliador, y debe homologar el acuerdo; según lo dicho en clase, el acuerdo del mediador no necesita homologación. \\
\midrule
\multicolumn{2}{p{0.92\textwidth}}{\textbf{Resumen:} el simulacro mostró cómo se abre la audiencia, cómo se fijan los hechos, cómo se depura la prueba, cómo se negocia y cómo se cierra un acuerdo que la jueza debe homologar. Dejó dos aprendizajes prácticos: la litis debe estar integrada con todas las partes, incluido el fiador, y el juez puede proponer fórmulas conciliatorias.} \\
\end{longtable}
\endgroup

% ==========================================================
% CAPÍTULO 6
% ==========================================================
\chapter[Evaluación y material pendiente]{Evaluación y material pendiente}

\section{Próxima clase y material pendiente}

La clase siguiente corresponde al \textbf{proceso ejecutivo} (ejecución o cobro de alquileres), otro proceso del que se dio una introducción al comienzo de la cursada. El examen final oral se fija para el 8 de octubre.

Quedó pendiente una clase sobre declaración de puro derecho, que no se dictó por coincidir con el día del estudiante. Respecto de la unidad nueve del cronograma, la docente menciona los siguientes contenidos: actitud del demandado, declaración de puro derecho, prueba, limitaciones, alegatos, abandono de la cosa, caducidad de instancia, sentencia de lanzamiento, recursos y resoluciones recurribles.

% TODO: contenido incompleto o ambiguo en las notas originales (la enumeración de la unidad nueve se leyó del cronograma de manera fragmentaria)

\section{Indicaciones para el examen final}

Para el examen se recomienda:

\begin{itemize}
    \item repasar el contenido de \textbf{procesal general};
    \item seguir las clases dictadas y los PowerPoint compartidos, que bastan para la preparación.
\end{itemize}

La forma del examen final es \textbf{oral}, con dos o tres estudiantes, ante los docentes de la cátedra. Se elige un número de la grilla (por ejemplo, del 1 al 50) y se pregunta sobre la unidad correspondiente.

\begin{exercise}{Pregunta de examen: la causa de puro derecho}
\textbf{Pregunta:} ¿En qué momento se puede declarar la causa de puro derecho en el proceso de desalojo? ¿Qué significa declararla de puro derecho?

\textbf{Respuesta expuesta en clase:} Hay dos momentos: el artículo 359, una vez trabada la litis desde el punto de vista sustancial, y la propia audiencia del artículo 360, según su inciso 6. Declarar la causa de puro derecho significa que el juez o jueza determinará cuál es el derecho aplicable a los hechos; ya no se discuten medios de prueba, porque no hay prueba por producir (fue producida como prueba anticipada, las partes desistieron de los medios o solo existe prueba documental). Si no hay contradicción, el juez dicta sentencia.
\end{exercise}

\section{Instancia de evaluación: coloquio}

Según lo expuesto por la docente, el reglamento obliga a ofrecer una segunda instancia de evaluación, que se organiza como \textbf{coloquio}. Se solicita al grupo que asista, dado que se trata de una forma de evaluación entre pares; quien rinde de forma individual queda desfasado. La materia continúa con procesal hasta diciembre.

% ---------------- CORNELL CAPÍTULO 6 ----------------
\section{Cuadro Cornell: evaluación}

\begingroup
\small
\begin{longtable}{
    >{\raggedright\arraybackslash}p{0.26\textwidth}
    >{\raggedright\arraybackslash}p{0.66\textwidth}
}
\toprule
\textbf{Preguntas / palabras clave} & \textbf{Notas / respuestas} \\
\midrule
\endfirsthead
\toprule
\textbf{Preguntas / palabras clave} & \textbf{Notas / respuestas} \\
\midrule
\endhead
\bottomrule
\endfoot
\textbf{¿Cómo es la evaluación?} &
Coloquio oral de dos o tres estudiantes, el 8 de octubre. Se elige un número de la grilla (1 a 50) y se responde sobre esa unidad. Conviene repasar procesal general y los PowerPoint de cada clase. \\[0.6em]

\textbf{¿En qué momentos se declara la causa de puro derecho?} &
En el art. 359, una vez trabada la litis desde el punto de vista sustancial, y en el art. 360, inciso 6, dentro de la audiencia. Declararla de puro derecho implica que se determina el derecho aplicable, porque no hay prueba por producir. \\[0.6em]

\textbf{¿Qué se verá en la próxima clase?} &
El proceso ejecutivo (ejecución o cobro de alquileres). \\
\midrule
\multicolumn{2}{p{0.92\textwidth}}{\textbf{Resumen:} la evaluación será un coloquio oral con preguntas sobre la unidad que salga de la grilla, como la de las dos oportunidades para declarar la causa de puro derecho (arts. 359 y 360, inciso 6). Se recomienda repasar procesal general y los PowerPoint de cada clase.} \\
\end{longtable}
\endgroup

\end{document}
